\subsection{A source-backed historical accounting check}
The February 2026 annual accounts allow a reconciliation that the October holdings snapshot does not. All amounts below are in \pounds million, as reported, not inferred from gilt purchase cost. The cash-flow statement gives
\[
13{,}304+134{,}833-101{,}372-25{,}051=21{,}714.
\]
Opening cash, net operating cash, loan repayments and interest paid thus reproduce closing cash exactly at the published precision. The carrying-value roll-forward is
\[
656{,}387-101{,}372+24{,}559-25{,}051=554{,}523,
\]
within \pounds1 million of the reported \pounds554,522 million closing loan, consistent with independently rounded statement lines. Accrued interest therefore matters even for the historical stock reconciliation. These are reproductions of the accounts, not new empirical discoveries.

For a narrowly defined sensitivity, hold the February principal and cash deposit fixed, let both reprice equally, and exclude all valuation, issuance, indemnity and behavioural responses. A permanent one-percentage-point change then changes the APF's direct annualised net financing charge by $0.01(L_p-C)$. The accounts report $L_p=\pounds553.2$ billion to one decimal place and $C=\pounds21{,}714$ million to the nearest million. Under conventional nearest-unit rounding their implied interval is
\[
0.01(L_p-C)\in[\pounds5.314355,\pounds5.315365]\ \text{billion per year}.
\]
This is an interval due to publication precision for a frozen-balance derivative, not a confidence interval, a forecast, the full public-sector net interest cost, or an October estimate. It demonstrates how a correctly defined accounting target can enter the bounding framework. Extending it over five years requires a cash/funding path and its shock response; the annual statements do not supply that path.

\paragraph{Independent source retrieval for the gilt ledger.} On 13 October an ordinary browser download from the official DMO report succeeded. Its internal heading is ``Gilts in issue on 12 October 2026''. All 105 ISINs, instrument types, redemption dates, nominal amounts and inflation-uplifted amounts agree with the supplied XML; the largest amount discrepancy is $10^{-11}$ million pounds, attributable to spreadsheet floating-point representation. The archived spreadsheet and row-by-row checker remove the earlier authentication gap for these computationally relevant gilt fields. This does not turn reconstructed bill timing into observed data or establish the APF funding model. The comparison is producer-coordinated source verification, not unaffiliated reproduction.
